% Part: model-theory
% Chapter: basics
% Section: reducts-and-expansions

\documentclass[../../../include/open-logic-section]{subfiles}

\begin{document}

\olfileid{mod}{bas}{red}
\section{المختزلات والتوسيعات}

قد تدعو الحاجة إلى مقابلة لغات تشترك في بعض الرموز، ثم مقابلة
ال!!{structure}s التي تؤولها. وأشهر ذلك أن تدخل رموز
!!a{language}~$\Lang{L}$ كلها في
!!a{language}~$\Lang{L'}$، فيكون $\Lang{L} \subseteq \Lang{L'}$.
وعندئذ نوسّع !!{structure}~$\Struct{M}$ للغة~$\Lang{L}$ إلى
!!{structure} للغة~$\Lang{L'}$، بأن نؤول الرموز الزائدة ونبقي
تأويل الرموز المشتركة على حاله. وبالعكس، نأخذ
!!{structure}~$\Struct{M'}$ للغة~$\Lang{L'}$ فنستخرج منها
!!{structure} للغة $\Lang{L}$ بإسقاط تأويل كل رمز ليس في~$\Lang{L}$؛
فكأننا «نسيناه».

\begin{defn}
\ollabel{defn:reduct}
لتكن $\Lang L \subseteq \Lang L'$، ولتكن $\Struct M$~!!{structure}
للغة~$\Lang L$، ولتكن $\Struct M'$~!!{structure} للغة~$\Lang L'$.
نسمي $\Struct M$ \emph{مختزَل}~$\Struct M'$ إلى~$\Lang L$، ونسمي
$\Struct M'$ \emph{توسيع}~$\Struct M$ إلى~$\Lang L'$، إذا وفقط إذا
\begin{enumerate}
\item $\Domain{{\OLId{latin-looped}{M}}} = \Domain{{\OLId{latin-looped}{M}}'}$.
\item لكل !!{constant}~${\OLId{latin-ordinary}{c}} \in \Lang L$، $\Assign{{\OLId{latin-ordinary}{c}}}{{\OLId{latin-looped}{M}}} =
  \Assign{{\OLId{latin-ordinary}{c}}}{{\OLId{latin-looped}{M}}'}$.
\item لكل !!{function}~${\OLId{latin-ordinary}{f}} \in \Lang L$، $\Assign{{\OLId{latin-ordinary}{f}}}{{\OLId{latin-looped}{M}}} =
  \Assign{{\OLId{latin-ordinary}{f}}}{{\OLId{latin-looped}{M}}'}$.
\item لكل !!{predicate}~${\OLId{latin-looped}{P}} \in \Lang L$، $\Assign{{\OLId{latin-looped}{P}}}{{\OLId{latin-looped}{M}}} =
  \Assign{{\OLId{latin-looped}{P}}}{{\OLId{latin-looped}{M}}'}$.
\end{enumerate}
\end{defn}

\begin{prop}
\ollabel{prop:reduct}
إذا كانت !!{structure}~$\Struct{M}$ للغة~$\Lang{L}$ مختزَل
!!{structure}~$\Struct{M'}$ للغة~$\Lang{L'}$، فإنه لجميع
ال!!{sentence}s~$!A$ في~$\Lang{L}$،
\[
\Sat{{\OLId{latin-looped}{M}}}{!A} \text{ إذا وفقط إذا كان } \Sat{{\OLId{latin-looped}{M}}'}{!A}.
\]
\end{prop}

\begin{proof}
  تمرين.
\end{proof}

\begin{prob}
أثبت \olref[mod][bas][red]{prop:reduct}.
\end{prob}

\begin{defn}
% تصحيح مصطلحي (C184-01): رُدّ n-place إلى «ذو/ذات n مواضع» بدل «n-محلية» المبهمة؛ والاختيار مفتوح لمراجعة خبير في مصطلح arity.
لنأخذ !!{structure}~$\Struct{M}$ للغة~$\Lang{L}$، ولتكن
$\Lang{L'} = \Lang{L} \cup \{{\OLId{latin-looped}{P}}\}$ توسيع~$\Lang{L}$ الناتج عن إضافة
!!{predicate} واحد ذو ${\OLId{latin-ordinary}{n}}$ مواضع~${\OLId{latin-looped}{P}}$، ولتكن ${\OLId{latin-looped}{R}} \subseteq \Domain{{\OLId{latin-looped}{M}}}^{\OLId{latin-ordinary}{n}}$
علاقة ذات ${\OLId{latin-ordinary}{n}}$ مواضع. عندئذ نرمز بـ $\Expan{{\OLId{latin-looped}{M}}}{{\OLId{latin-looped}{R}}}$ إلى التوسيع~$\Struct{M'}$
لـ~$\Struct{M}$ الذي فيه $\Assign{{\OLId{latin-looped}{P}}}{{\OLId{latin-looped}{M}}'} = {\OLId{latin-looped}{R}}$.
\end{defn}

\end{document}
